% GENERATED FILE. Edit canonical lessons, then run npm run generate:books.
% canonical-source-hash: fnv1a64:bce2d2bbca31faff
% canonical-lessons: ML-C441-sadhyamaya, ML-C441-asadhyamaya, ML-C441-sadhyamakkuka, ML-C441-annaneyenkil, ML-C441-mikkavarum

\chapter{Possible or Not}
\label{ch:ml-possible-or-not}
\hlchaptermodality{441}

\begin{chapteropening}
\textbf{By the end of this chapter:} \emph{I can say possible, impossible, to make possible, in that case and most probably.}

Complete the last lesson of chapter 441: I can say possible, impossible, to make possible, in that case and most probably.
\end{chapteropening}

\section[\texorpdfstring{\ml{സാധ്യമായ} (sādhyamāya)}{സാധ്യമായ (sādhyamāya)}]{\ml{സാധ്യമായ} (sādhyamāya) — possible}

\label{lesson:ML-C441-sadhyamaya}



\begin{quote}
Before the new one: say the Malayalam for if not, then the Malayalam for or else.
\end{quote}

\subsection*{You'll want to know: \ml{സാധ്യമായ}}
\textbf{\ml{സാധ്യമായ}} — \emph{sādhyamāya} — \textquotedblleft{}possible\textquotedblright{}.

\begin{quote}
\ml{സാധ്യമായ} \ml{എല്ലാ} \ml{വഴികളും} \ml{അവർ} \ml{നോക്കി}.
\end{quote}

\begin{quote}
\emph{sādhyamāya ellā vaḻikaḷuṁ avar nōkki.}
\end{quote}

\begin{quote}
They tried every possible way.
\end{quote}

As a predicate: \ml{അത്} \ml{സാധ്യമാണ്}, that is possible.

\begin{cousinweb}
\textbf{\ml{സാധ്യമായ}} (\emph{sādhyamāya}) is built from Sanskrit \emph{sādhya}, \textquotedblleft{}to be mastered, to be accomplished\textquotedblright{} (Monier-Williams), with the adjective ending -\ml{ആയ}.
\end{cousinweb}

\begin{grammarlens}[title={What you've built}]
One more word for saying what happens if.
\end{grammarlens}

\subsection*{Your turn}
\begin{itemize}
\raggedright
  \item \emph{Say it:} \emph{sādhyamāya}
  \item \emph{Say it:} \emph{sādhyamāya}, once more
  \item \emph{Say it:} \emph{sādhyamāya ellā vaḻikaḷuṁ avar nōkki.} — They tried every possible way.
  \item \emph{Recall:} say the Malayalam for if not, then the Malayalam for or else, then say \emph{sādhyamāya} again
\end{itemize}

\begin{tcolorbox}[breakable,colback=teal!4,colframe=teal!35!black,title={Before you move on}]
What does \ml{സാധ്യമായ} mean? (\textquotedblleft{}Possible\textquotedblright{}.) Say it once more.
\end{tcolorbox}
\section[\texorpdfstring{\ml{അസാധ്യമായ} (asādhyamāya)}{അസാധ്യമായ (asādhyamāya)}]{\ml{അസാധ്യമായ} (asādhyamāya) — impossible}

\label{lesson:ML-C441-asadhyamaya}



\begin{quote}
Before the new one: say the Malayalam for or else, then the Malayalam for possible.
\end{quote}

\subsection*{You'll want to know: \ml{അസാധ്യമായ}}
\textbf{\ml{അസാധ്യമായ}} — \emph{asādhyamāya} — \textquotedblleft{}impossible\textquotedblright{}.

\begin{quote}
\ml{അവന്} \ml{അസാധ്യമായ} \ml{ഒന്നുമില്ല}.
\end{quote}

\begin{quote}
\emph{avanŭ asādhyamāya onnumilla.}
\end{quote}

\begin{quote}
Nothing is impossible for him.
\end{quote}

As a predicate: \ml{അത്} \ml{അസാധ്യമാണ്}.

\begin{cousinweb}
\textbf{\ml{അസാധ്യമായ}} (\emph{asādhyamāya}) is built from Sanskrit \emph{asādhya}, \textquotedblleft{}not to be accomplished\textquotedblright{} (Monier-Williams), with the adjective ending -\ml{ആയ}.
\end{cousinweb}

\begin{grammarlens}[title={What you've built}]
One more word for saying what happens if.
\end{grammarlens}

\subsection*{Your turn}
\begin{itemize}
\raggedright
  \item \emph{Say it:} \emph{asādhyamāya}
  \item \emph{Say it:} \emph{asādhyamāya}, once more
  \item \emph{Say it:} \emph{avanŭ asādhyamāya onnumilla.} — Nothing is impossible for him.
  \item \emph{Recall:} say the Malayalam for or else, then the Malayalam for possible, then say \emph{asādhyamāya} again
\end{itemize}

\begin{tcolorbox}[breakable,colback=teal!4,colframe=teal!35!black,title={Before you move on}]
What does \ml{അസാധ്യമായ} mean? (\textquotedblleft{}Impossible\textquotedblright{}.) Say it once more.
\end{tcolorbox}
\section[\texorpdfstring{\ml{സാധ്യമാക്കുക} (sādhyamākkuka)}{സാധ്യമാക്കുക (sādhyamākkuka)}]{\ml{സാധ്യമാക്കുക} (sādhyamākkuka) — to make possible}

\label{lesson:ML-C441-sadhyamakkuka}



\begin{quote}
Before the new one: say the Malayalam for possible, then the Malayalam for impossible.
\end{quote}

\subsection*{You'll want to know: \ml{സാധ്യമാക്കുക}}
\textbf{\ml{സാധ്യമാക്കുക}} — \emph{sādhyamākkuka} — \textquotedblleft{}to make possible\textquotedblright{}.

\begin{quote}
\ml{നാട്ടുകാരുടെ} \ml{സഹായം} \ml{ആ} \ml{പാലം} \ml{സാധ്യമാക്കി}.
\end{quote}

\begin{quote}
\emph{nāṭṭukāruṭe sahāyaṁ ā pālaṁ sādhyamākki.}
\end{quote}

\begin{quote}
The local people's help made that bridge possible.
\end{quote}

Past: \ml{സാധ്യമാക്കി}, \emph{sādhyamākki}.

\ml{ആക്കുക}, to make, turns an adjective into a verb.

\begin{cousinweb}
\textbf{\ml{സാധ്യമാക്കുക}} (\emph{sādhyamākkuka}) is built from Sanskrit \emph{sādhya}, \textquotedblleft{}to be accomplished\textquotedblright{} (Monier-Williams), with \ml{ആക്കുക}, \textquotedblleft{}to make\textquotedblright{}, native Dravidian (DEDR 333).
\end{cousinweb}

\begin{grammarlens}[title={What you've built}]
One more word for saying what happens if.
\end{grammarlens}

\subsection*{Your turn}
\begin{itemize}
\raggedright
  \item \emph{Say it:} \emph{sādhyamākkuka}
  \item \emph{Say it:} \emph{sādhyamākkuka}, once more
  \item \emph{Say it:} \emph{nāṭṭukāruṭe sahāyaṁ ā pālaṁ sādhyamākki.} — The local people's help made that bridge possible.
  \item \emph{Recall:} say the Malayalam for possible, then the Malayalam for impossible, then say \emph{sādhyamākkuka} again
\end{itemize}

\begin{tcolorbox}[breakable,colback=teal!4,colframe=teal!35!black,title={Before you move on}]
What does \ml{സാധ്യമാക്കുക} mean? (\textquotedblleft{}To make possible\textquotedblright{}.) Say it once more.
\end{tcolorbox}
\section[\texorpdfstring{\ml{അങ്ങനെയെങ്കിൽ} (aṅṅaneyeṅkil)}{അങ്ങനെയെങ്കിൽ (aṅṅaneyeṅkil)}]{\ml{അങ്ങനെയെങ്കിൽ} (aṅṅaneyeṅkil) — in that case}

\label{lesson:ML-C441-annaneyenkil}



\begin{quote}
Before the new one: say the Malayalam for impossible, then the Malayalam for to make possible.
\end{quote}

\subsection*{You'll want to know: \ml{അങ്ങനെയെങ്കിൽ}}
\textbf{\ml{അങ്ങനെയെങ്കിൽ}} — \emph{aṅṅaneyeṅkil} — \textquotedblleft{}in that case\textquotedblright{}.

\begin{quote}
\ml{അവൻ} \ml{വരില്ലെന്ന്} \ml{പറഞ്ഞു}, \ml{അങ്ങനെയെങ്കിൽ} \ml{നമ്മൾ} \ml{നേരത്തെ} \ml{പോകും}.
\end{quote}

\begin{quote}
\emph{avan varillennŭ paṟaññu, aṅṅaneyeṅkil nammaḷ nēratte pōkuṁ.}
\end{quote}

\begin{quote}
He said he won't come; in that case we will leave early.
\end{quote}

\ml{അങ്ങനെ}, like that, with -\ml{എങ്കിൽ}, if: if it is like that.

\begin{cousinweb}
\textbf{\ml{അങ്ങനെയെങ്കിൽ}} (\emph{aṅṅaneyeṅkil}) is native Dravidian: \ml{അങ്ങനെ}, \textquotedblleft{}in that way\textquotedblright{}, from the far demonstrative \ml{അ}- (DEDR 1), with -\ml{എങ്കിൽ}, \textquotedblleft{}if\textquotedblright{}.
\end{cousinweb}

\begin{grammarlens}[title={What you've built}]
One more word for saying what happens if.
\end{grammarlens}

\subsection*{Your turn}
\begin{itemize}
\raggedright
  \item \emph{Say it:} \emph{aṅṅaneyeṅkil}
  \item \emph{Say it:} \emph{aṅṅaneyeṅkil}, once more
  \item \emph{Say it:} \emph{avan varillennŭ paṟaññu, aṅṅaneyeṅkil nammaḷ nēratte pōkuṁ.} — He said he won't come; in that case we will leave early.
  \item \emph{Recall:} say the Malayalam for impossible, then the Malayalam for to make possible, then say \emph{aṅṅaneyeṅkil} again
\end{itemize}

\begin{tcolorbox}[breakable,colback=teal!4,colframe=teal!35!black,title={Before you move on}]
What does \ml{അങ്ങനെയെങ്കിൽ} mean? (\textquotedblleft{}In that case\textquotedblright{}.) Say it once more.
\end{tcolorbox}
\section[\texorpdfstring{\ml{മിക്കവാറും} (mikkavāṟuṁ)}{മിക്കവാറും (mikkavāṟuṁ)}]{\ml{മിക്കവാറും} (mikkavāṟuṁ) — most probably; mostly}

\label{lesson:ML-C441-mikkavarum}



\begin{quote}
Before the new one: say the Malayalam for to make possible, then the Malayalam for in that case.
\end{quote}

\subsection*{You'll want to know: \ml{മിക്കവാറും}}
\textbf{\ml{മിക്കവാറും}} — \emph{mikkavāṟuṁ} — \textquotedblleft{}most probably; mostly\textquotedblright{}.

\begin{quote}
\ml{മിക്കവാറും} \ml{നാളെ} \ml{മഴ} \ml{പെയ്യും}.
\end{quote}

\begin{quote}
\emph{mikkavāṟuṁ nāḷe maḻa peyyuṁ.}
\end{quote}

\begin{quote}
It will most probably rain tomorrow.
\end{quote}

With the future -\ml{ഉം} it means most probably; with numbers, nearly all.

\begin{cousinweb}
\textbf{\ml{മിക്കവാറും}} (\emph{mikkavāṟuṁ}) is native Dravidian, built on \ml{മിക്ക}, \textquotedblleft{}the greater part\textquotedblright{}, from \ml{മികുക}, \textquotedblleft{}to exceed\textquotedblright{} (DEDR 4838): for the greater part.
\end{cousinweb}

\begin{grammarlens}[title={What you've built}]
One more word for saying what happens if.
\end{grammarlens}

\subsection*{Your turn}
\begin{itemize}
\raggedright
  \item \emph{Say it:} \emph{mikkavāṟuṁ}
  \item \emph{Say it:} \emph{mikkavāṟuṁ}, once more
  \item \emph{Say it:} \emph{mikkavāṟuṁ nāḷe maḻa peyyuṁ.} — It will most probably rain tomorrow.
  \item \emph{Recall:} say the Malayalam for to make possible, then the Malayalam for in that case, then say \emph{mikkavāṟuṁ} again
\end{itemize}

\begin{tcolorbox}[breakable,colback=teal!4,colframe=teal!35!black,title={Before you move on}]
What does \ml{മിക്കവാറും} mean? (\textquotedblleft{}Most probably; mostly\textquotedblright{}.) Say it once more.
\end{tcolorbox}
